\section{部署核对清单}
真正把课程里的分析落到服务侧，团队通常要在上线前回答几个问题：当前业务的 prompt 长短分布是什么？主要优化 TTFT 还是 TPS？KV cache 是否会在高并发时触发显存抖动？这些问题如果不提前梳理，往往会让系统在真实流量下暴露出与 benchmark 完全不同的瓶颈。

\begin{importantbox}{上线前至少核对四件事}
\begin{itemize}
  \item 用户请求以长 prompt 还是长回答为主
  \item 现有硬件瓶颈更偏向算力还是带宽
  \item KV cache 和 batch 策略是否匹配目标并发
  \item 监控面板是否同时覆盖 TTFT、TPS 与显存占用
\end{itemize}
\end{importantbox}

\subsection{本章小结}
Prefill/Decode 的理论差异最终都会回到部署清单上。只有把工作负载、硬件和指标统一到一个核对流程里，推理优化才真正可执行。

\newpage
